\begin{abstract}
Benchmarks for molecular property prediction are scored against experimental labels, yet the same compound measured in different laboratories routinely disagrees. We ask how much of the difference between models lies inside this label noise, and whether pretrained molecular embeddings still beat simple ECFP-based baselines once scores are expressed relative to the measurement limit. We estimate a \emph{noise ceiling} per endpoint from replicate measurements of the same molecules reported in \emph{independent ChEMBL documents}, rather than assuming an error size, and convert model scores to a \emph{fraction of the ceiling}. We release a replicate-based dataset of molecules not present in existing public benchmarks, a benchmark of small models (ECFP with GBM/RF/MLP, GNN) and frozen foundation-model embeddings with a trained head under random, scaffold and temporal splits, and a reporting protocol with minimum detectable effects. Results: \tbd{} (no results have been computed yet; every number in this draft is a placeholder).
\end{abstract}
